\documentclass[letterpaper]{article} 
\usepackage[utf8]{inputenc}
\linespread{0.85}
\usepackage[T1]{fontenc}
\usepackage{amsmath}
\usepackage{amsfonts}
\usepackage{amssymb}
\usepackage{array}
\usepackage{fontspec}
\usepackage{booktabs}
\usepackage{hyperref}
\usepackage[version=4]{mhchem}
\usepackage{stmaryrd}
\usepackage[dvipsnames]{xcolor}
\colorlet{LightRubineRed}{RubineRed!70}
\colorlet{Mycolor1}{green!10!orange}
\definecolor{Mycolor2}{HTML}{00F9DE}
\usepackage{graphicx}
\usepackage{amsmath}
\usepackage{graphicx}
\usepackage{capt-of}
\usepackage{lipsum}
\usepackage{fancyvrb}
\usepackage{tabularx}
\usepackage{listings}
\usepackage[export]{adjustbox}
\graphicspath{ {./images/} }
\usepackage[utf8]{inputenc}
\usepackage[english]{babel}
\usepackage{float}
\usepackage{lipsum}
\usepackage{graphicx}
\usepackage{float}
\usepackage[margin=0.7in]{geometry}
\usepackage{amsmath}
\usepackage{graphicx}
\usepackage{capt-of}
\usepackage{tcolorbox}
\usepackage{lipsum}
\usepackage{graphicx}
\usepackage{float}
\usepackage{listings}
\definecolor{deepred}{RGB}{139,0,0}
\usepackage{hyperref} 
\usepackage{xcolor} % For custom colors
\lstset{
	language=Python,                % Choose the language (e.g., Python, C, R)
	basicstyle=\ttfamily\small, % Font size and type
	keywordstyle=\color{blue},  % Keywords color
	commentstyle=\color{gray},  % Comments color
	stringstyle=\color{red},    % String color
	numbers=left,               % Line numbers
	numberstyle=\tiny\color{gray}, % Line number style
	stepnumber=1,               % Numbering step
	breaklines=true,            % Auto line break
	backgroundcolor=\color{black!5}, % Light gray background
	frame=single,               % Frame around the code
}
\usepackage{float}
\usepackage[]{amsthm} %lets us use \begin{proof}
	\usepackage[]{amssymb} %gives us the character \varnothing

	\title{Homework 5, MATH 4155}
	\author{Zongyi Liu}
	\date{Tue, Feb 24, 2026}
	\begin{document}
		\maketitle
		\section{Question 1}
		
		Consider the coin-toss model with success probability $p \in (0,1)$, that is, take $\Omega = \{0,1\}^n$ as your sample space, take $\mathcal{F} = \mathcal{P}(\Omega)$, set $X_i(\omega) = \omega_i$ for $i=1,\cdots,n$ and construct a probability measure $\mathbb{P}$ such that
		\[
		\mathbb{P}\bigl(X_1 = x_1, \cdots, X_n = x_n\bigr)
		=
		\left.\bigl(p^{s}(1-p)^{\,n-s}\bigr)\right|_{s=\sum_{i=1}^n x_i},
		\qquad
		\forall (x_1,\cdots,x_n)\in\{0,1\}^n.
		\]
		
		As we saw in class, under $\mathbb{P}$ the random variables $X_1,\cdots,X_n$ are independent Bernoulli random variables. Set $S_n(\omega):=\sum_{i=1}^n X_i(\omega)$ and fix an integer $0\le k\le n$.
		
		\begin{enumerate}
			\item[(i)] Find the conditional distribution of $X_1,\cdots,X_n$ given that $k$ successes occur among the $n$ tosses. In other words, compute
			\[
			\mathbb{P}\bigl(X_1=x_1,\cdots,X_n=x_n \mid S_n=k\bigr),
			\qquad
			(x_1,\cdots,x_n)\in\{0,1\}^n.
			\]
			
			\item[(ii)] Compute $\mathbb{E}\bigl(X_i \mid S_n=k\bigr)$, $i=1,\cdots,n$.
			
			\item[(iii)] For any given $m\in\{1,\cdots,n-1\}$ and $(x_1,\cdots,x_m)\in\{0,1\}^m$, compute the conditional probabilities
			\[
			\mathbb{P}\bigl(X_{m+1}=1 \mid X_1=x_1,\cdots,X_m=x_m,S_n=k\bigr).
			\]
		\end{enumerate}
		
		\textbf{Answer}
		
		\clearpage
		
		\section{Question 2}
		
		Rejection Sampling. Let $X$ and $Y$ be random variables with possible values in the countable set $\mathfrak{Z}=\left\{z_{1}, z_{2}, \cdots\right\}$ and such that
		
		$$
		p(z):=\mathbb{P}(X=z), \quad q(z):=\mathbb{P}(Y=z) \quad \text { satisfy } \quad p(z) \leq c q(z), \quad \forall z \in \mathfrak{Z}
		$$
		
		for some real constant $c>0$. Let $U$ be independent of $Y$ with uniform distribution on $[0,1]$, and set $\lambda(z):=p(z) /(c q(z)), z \in \mathfrak{Z}$.
		
		Show that the conditional distribution of the random variable $Y$ given the event $\{U \leq \lambda(Y)\}$, is the same as the distribution of the random variable $X$; to wit,
		
		$$
		\mathbb{P}(Y=z \mid U \leq \lambda(Y))=\mathbb{P}(X=z), \quad z \in \mathfrak{Z} .
		$$
		
		
		\textbf{Answer}
		
		\clearpage
		
		\section{Question 3}
		
		Use the Tonelli theorem to show that, for any random variable $X: \Omega \rightarrow [0, \infty)$,
		
		$$
		\mathbb{E}(X)=\int_{0}^{\infty} \mathbb{P}(X>t) \mathrm{d} t
		$$
		
		or, more generally for any $r \in(0, \infty)$ :
		
		$$
		\mathbb{E}\left(X^{r}\right)=r \int_{0}^{\infty} t^{r-1} \mathbb{P}(X>t) \mathrm{d} t
		$$
		
		Use this expression to compute - painlessly - the moments of the exponential distribution. (\emph{Hint}: Write $X(\omega)=\int_{0}^{\infty} 1_{[0, X(\omega))} d t$.)
		
		
		\textbf{Answer}
		
		\clearpage
		
		\section{Question 4}
		
		(*) For each $n \in \mathbb{N}$, we set
		
		$$
		A_{n}=\left\{\frac{m}{n}: m \in \mathbb{N}\right\}
		$$
		
		Determine the sets
		
		$$
		\left\{A_{n}, \text { ev. }\right\}:=\bigcup_{k \in \mathbb{N}} \bigcap_{n \geq k} A_{n}, \quad\left\{A_{n}, \text { i.o. }\right\}:=\bigcap_{k \in \mathbb{N}} \bigcup_{n \geq k} A_{n}
		$$
		
		("eventually" and "infinitely often", respectively).
		
		\textbf{Answer}
		
		\clearpage
		
		\section{Question 5}
		
		(*) Let $X_{1}, X_{2}, \cdots$ be random variables with values in [ $0, \infty$ ) and known distribution functions $F_{n}(\xi)=\mathbb{P}\left(X_{n} \leq \xi\right), \xi \in \mathbb{R}$. Find constants $c_{1}>0, c_{2}>0, \cdots$ such that
		
		$$
		\mathbb{P}\left(\lim _{n \rightarrow \infty}\left(X_{n} / c_{n}\right) \quad \text { exists, and equals } 0\right)=1 .
		$$
		
		(\emph{Hint}: Recall the Borel-Cantelli Lemma; provide a very careful argument.)
		
		\textbf{Answer}
		
		\clearpage
		
		\section{Question 6}
		
		{\color{deepred}\fontspec{Georgia} (*) Lebesgue Measure Cannot be Extended to the Power Set $\mathcal{P}([0,1))$ under the Axiom of Choice}
		
		Let $\mathbb{P}$ denote Lebesgue measure on $([0,1), \mathcal{B}([0,1)))$. For any given $a \in \mathbb{R}$, define the ``shift mapping'' $T_a(x) = x + a \bmod 1$, for $x \in [0,1)$; for instance, $T_{0.9}(0.3) = 0.2$.
		
		\begin{enumerate}
			\item[(i)] Show that each $T_a$ is measurable with respect to $\mathcal{B}([0,1))$, and that Lebesgue measure is invariant under $T_a$, that is, $\mathbb{P} \circ T_a^{-1} = \mathbb{P}$.
			
			\item[(ii)] Now let $\mathcal{F}$ be any $\sigma$-algebra of subsets of $[0,1)$ for which $T_a$ is measurable, and assume $\mathbb{P}$ is defined and shift-invariant (as in (i) above) on all of $\mathcal{F}$.
			
			Show then that any set $A \in \mathcal{F}$ with $\mathbb{P}(A) > 0$ contains points $x, y$ such that $x \neq y$ and $x - y$ is rational.
			
			(Hint: Show that $\{T_q^{-1}(A) : q \in \mathbf{Q} \cap [0,1)\}$ is a family of disjoint sets, and use part (i).)
			
			\item[(iii)] Conclude that Lebesgue measure cannot be extended to $\mathcal{P}([0,1))$, the collection of all subsets of the unit interval $[0,1)$, all the while retaining its characteristic property of shift-invariance.
			
			(\emph{Hint}: Consider the equivalence relation
			\[
			x \sim y \quad \Longleftrightarrow \quad x - y \text{ is rational},
			\]
			and invoke the Axiom of Choice${}^1$ to obtain a set $A \in \mathcal{P}([0,1))$ of representatives from each of the induced equivalence classes. This gives a decomposition
			\[
			[0,1) = \bigcup_{r \in \mathbf{Q} \cap [0,1)} T_r^{-1}(A).
			\]
			What implications for $\mathbb{P}(A)$ does this decomposition have?)
		\end{enumerate}
		
		\textbf{Answer}
		
		\clearpage
		
		\section{Question 7}
		
		{\color{deepred}\fontspec{Georgia} (*) Growth of Exponential Random Variables, or The Unreasonable Effectiveness of Borel-Cantelli}
		
		 Suppose $X_{1}, X_{2}, \cdots$ are independent random variables with common exponential distribution $\mathbb{P}\left(X_{n}>\xi\right)=e^{-\xi}, \xi \geq 0$; let us say, the durations of successive telephone conversations, for which the exponential is a good model.
		
		Let us look at the maximum $M_{n}=\max _{1 \leq k \leq n} X_{k}$ of the first $n$ of these variables. It should be clear that this random variable $M_{n}$ increases to infinity a.e., as $n \rightarrow \infty$ (give a brief argument).
		
		But at what rate? Show that
		
		$$
		\lim _{n \rightarrow \infty}\left(\frac{M_{n}}{\log n}\right)=1, \quad \text { a.e. }
		$$
		
		which provides the answer to this question: the quantity $M_{n}$ increases to infinity, but pretty slowly. (\emph{Hint}: Use Borel-Cantelli, to establish first
		
		$$
		\limsup _{n \rightarrow \infty}\left(\frac{X_{n}}{\log n}\right)=1
		$$
		
		In fact, show that for every $\varepsilon \in(0,1)$ we have
		
		$$
		\begin{array}{ll}
			\mathbb{P}\left(X_{n}>(1-\varepsilon) \log n,\right. & \text { i.o. })=1 \\
			\mathbb{P}\left(X_{n}>(1+\varepsilon) \log n,\right. & \text { i.o. })=0 \\
			\mathbb{P}\left(M_{n}<(1-\varepsilon) \log n,\right. & \text { i.o. })=0
		\end{array}
		$$
		
		Deduce from all this the double a.e. inequality
		
		$$
		\limsup _{n \rightarrow \infty}\left(\frac{M_{n}}{\log n}\right) \leq 1 \leq \liminf _{n \rightarrow \infty}\left(\frac{M_{n}}{\log n}\right),
		$$
		
		and argue that this settles the claim.)
		
		\textbf{Answer}
		
		\clearpage
		
		\section{Question 8}
		
		{\color{deepred}\fontspec{Georgia} Typed Notes 3.9} 
		
		If $N$ is a \textsc{Poisson} process with parameter $\lambda>0$, and if $0\le s<t<\infty$ are given real numbers, compute
		\[
		\mathbb{P}\bigl(N(s)=1,\; N(t)=2\bigr).
		\]
		
	
		
		\textbf{Answer}
		
		\clearpage
		
		\section{Question 9}
		
			{\color{deepred}\fontspec{Georgia} Typed Notes 3.10} 
			
			Show that, if $N$ is a \textsc{Poisson} process with parameter $\lambda>0$, and if $t>0$ is a given real number, we have
		\[
		\mathbb{P}\bigl(N(t)=n\bigr)=\mathbb{P}\bigl(S_n\le t<S_{n+1}\bigr)
		= e^{-\lambda t}\frac{(\lambda t)^n}{n!},\qquad k=0,1,\cdots .
		\]
		
		(\textit{Hint:} Argue by induction the validity of the representation
		\[
		\mathbb{P}(S_n>t)=\int_t^\infty \lambda e^{-\lambda s}\frac{(\lambda s)^{n-1}}{(n-1)!}\,ds
		= e^{-\lambda t}\sum_{k=0}^{n-1}\frac{(\lambda t)^k}{k!},\qquad n=1,2,\cdots
		\]
		for the Gamma $\Gamma(\lambda,n)$ distribution in (3.17).)
		
		\textbf{Answer}
		
		\clearpage
		
		\section{Question 10}
		
		{\color{deepred}\fontspec{Georgia} Walsh 5.23} 
		
		Let $X_1, X_2, \ldots$ be a sequence of i.i.d.\ random variables with $E\{|X_1|\}=\infty$.
		\begin{enumerate}
			\item[(a)] Show that $\displaystyle \limsup_{n\to\infty}\frac{|X_n|}{n}=\infty$ a.s.
			\item[(b)] Deduce that $\displaystyle \limsup_{n\to\infty}\frac{|X_1+\cdots+X_n|}{n}=\infty$.
		\end{enumerate}
		
		\textbf{Answer}
		
		\clearpage
		
	\end{document}
